\documentclass[12pt]{article}
\usepackage{amsmath,amssymb,amsfonts}
\usepackage[margin=1in]{geometry}

\begin{document}

\textbf{Chapter 6, Problem 35.} The $\theta$-function (or step function) is defined as
\[
\theta(x) = \begin{cases} 1 & \text{for } x > 0, \\ \frac{1}{2} & \text{for } x = 0, \\ 0 & \text{for } x < 0. \end{cases}
\]

Prove:
\begin{itemize}
\item[(a)] $\theta(x) = \frac{1}{2} + \frac{1}{2\pi i}\lim_{\epsilon\to 0^+}\int_{-\infty}^{\infty}\frac{e^{ikx}}{k-i\epsilon}\,dk$
\item[(b)] $\frac{d}{dx}\theta(x) = \delta(x)$
\end{itemize}

[\textit{Hint:} See Eq.\ (5.63).]

\bigskip
\textbf{Solution.}

\textbf{(a)} Consider the integral
\[
I = \int_{-\infty}^{\infty}\frac{e^{ikx}}{k - i\epsilon}\,dk, \quad \epsilon > 0.
\]

The pole is at $k = i\epsilon$ (in the upper half $k$-plane).

\textbf{Case $x > 0$:} Close the contour in the upper half-plane (where $e^{ikx}$ decays for $\text{Im}(k) > 0$). The pole $k = i\epsilon$ is enclosed. By the residue theorem:
\[
I = 2\pi i \cdot e^{i(i\epsilon)x} = 2\pi i\,e^{-\epsilon x}.
\]
As $\epsilon \to 0^+$: $I \to 2\pi i$. So $\frac{1}{2\pi i}I = 1$, and $\frac{1}{2}+\frac{1}{2\pi i}I \to \frac{1}{2}+1$...

Wait, let me reconsider. We need $\theta(x) = \frac{1}{2} + \frac{1}{2\pi i}\lim I$. For $x > 0$: this gives $\frac{1}{2} + 1 = \frac{3}{2}$, which is wrong.

Let me reconsider the representation. The standard result is:
\[
\theta(x) = \frac{1}{2\pi i}\lim_{\epsilon\to 0^+}\int_{-\infty}^{\infty}\frac{e^{ikx}}{k - i\epsilon}\,dk.
\]

For $x > 0$: close in upper half-plane, pick up pole at $k = i\epsilon$: $I = 2\pi i \cdot e^{-\epsilon x} \to 2\pi i$. So $\theta(x) = 1$. $\checkmark$

For $x < 0$: close in lower half-plane (where $e^{ikx}$ decays). No poles enclosed. $I = 0$, so $\theta(x) = 0$. $\checkmark$

For $x = 0$: $I = \int_{-\infty}^{\infty}\frac{dk}{k-i\epsilon}$. Using a principal value argument and a semicircular contour, this gives $I = \pi i$, so $\theta(0) = 1/2$. $\checkmark$

Therefore the correct formula is:
\[
\theta(x) = \frac{1}{2\pi i}\lim_{\epsilon\to 0^+}\int_{-\infty}^{\infty}\frac{e^{ikx}}{k-i\epsilon}\,dk. \quad \checkmark
\]

(The problem's formula with the $\frac{1}{2}+$ may involve a different convention or a different integral representation.)

\textbf{(b)} To show $\theta'(x) = \delta(x)$:

From the Fourier representation of $\delta$:
\[
\delta(x) = \frac{1}{2\pi}\int_{-\infty}^{\infty}e^{ikx}\,dk.
\]

Differentiate $\theta(x)$:
\[
\frac{d}{dx}\theta(x) = \frac{1}{2\pi i}\lim_{\epsilon\to 0^+}\int_{-\infty}^{\infty}\frac{ik\,e^{ikx}}{k-i\epsilon}\,dk.
\]

As $\epsilon \to 0$:
\[
\frac{ik}{k-i\epsilon} \to \frac{ik}{k} = i \quad \text{(distributional sense)}.
\]

So:
\[
\theta'(x) = \frac{1}{2\pi i}\int_{-\infty}^{\infty}ie^{ikx}\,dk = \frac{1}{2\pi}\int_{-\infty}^{\infty}e^{ikx}\,dk = \delta(x). \quad
\]

\end{document}
